\begin{abstract}
Power line communication (PLC) uses existing electrical conductors as a medium for data transmission. This makes the technology attractive in situations where installing separate communication wiring would be expensive or impractical. However, power lines are not designed as communication channels, which creates several challenges, including noise, changing line impedance, signal attenuation, coupling safety, and electromagnetic emission limits.

This paper gives an overview of PLC by examining these challenges and the methods used to work around them. First, the basic operating principle of PLC is introduced. Then, typical difficulties of using power lines for communication are discussed. The paper also considers several application areas, including home PLC implementations, smart metering, and industrial examples such as feedback communication over motor PWM wiring. These examples are used to show how PLC remains useful as a niche technology when the existing power infrastructure can be reused effectively.
\end{abstract}